%% main.tex — Anvil paper (thread: conformal-antenna-diffopt, version .1)
%% Drafted by anvil:paper-draft. Numbers trace to the committed artifact
%% benchmarks/patch_antenna_conformal/conformal_results.toml (branch
%% feature/issue-650, commit 524db3b) and crates/geode-core source.

\documentclass{anvil-paper}

\title{Differentiable-by-Construction FEM for Freeform Open-Radiator Design:
A Curved Conformal Antenna Structured-Grid Inverse Design Cannot Reach}
\author{Robb Walters}
\date{\today}

\begin{document}
\maketitle

\begin{abstract}
Differentiable electromagnetic (EM) simulators have made photonic inverse
design routine, but the dominant tools optimize a permittivity \emph{density}
field on a structured Yee/Cartesian grid, a representation that staircases
curved metal and is a poor fit for open radiators terminated by a perfectly
matched layer (PML). We present a tensor-compiled, differentiable-by-construction
3-D N\'ed\'elec H(curl) finite-element method (FEM) that exposes an exact,
single-factorization discrete \emph{shape} (moving-boundary / node-motion)
adjoint through the full open-radiator forward: a matched box-UPML tensor
material, lossy complex-$\varepsilon$ FR-4, and a lumped feed port assembled into
one complex-symmetric pencil $A(\omega)=K(\nu)-\omega^2 M(\varepsilon)+(j\omega/Z_s)S_p$.
Forward-mode ``Dual-twin'' element kernels give exact
$\partial K/\partial X,\ \partial M/\partial X,\ \partial b/\partial X$, so the
transpose adjoint reuses the single forward factorization
($n_{\text{fact}}=1$ per band frequency). We drive a high-DOF harmonic
mesh-morph parametrization (73 active boundary DOFs) to gradient-design a
\emph{freeform curved-metal conformal} radiator, a geometry class density-based
differentiable EM represents poorly. The optimizer reaches a genuine
target-hit terminal that drives the entire design band below $-10$~dB return
loss (worst-of-band $-5.51 \to -12.06$~dB), with the band gradient FD-validated
at relative error $1.17\times10^{-10}$, enforced passivity $|S_{11}|\le1$, and a
non-inversion guard. We position the narrow, honestly-scoped novelty against
prior EM shape adjoints and the structured-grid density ecosystem.
\end{abstract}

\section{Introduction}
\label{sec:intro}

Inverse design has become the default workflow for nanophotonic components:
given a differentiable forward Maxwell solver, one back-propagates a
figure-of-merit to a large parameter vector and lets gradient descent discover
non-intuitive structures. The canonical ``killer'' demonstrations --- the
compact wavelength demultiplexer of Piggott et al.~\cite{piggott2015} and the
high-NA achromatic metalens of Chung and Miller~\cite{chung2020} --- are
guided-wave or periodic \emph{photonic} problems, and the workhorse tools
(\textsc{ceviche}, Meep's adjoint solver, SPINS-B~\cite{su_spinsb},
FDTDX~\cite{fdtdx}, invrs-gym~\cite{invrsgym}) optimize a permittivity
\emph{density} field discretized on a structured Yee/Cartesian grid. This
representation is a superb fit for topology optimization of dielectric
photonics, but it has two structural mismatches with a large and practically
important class of RF/antenna problems: (i) it \emph{staircases} curved metal,
so a conformal radiator wrapped around a cylinder is only ever approximated at a
density-limited grid cost; and (ii) it optimizes ``where is there material''
rather than ``where is the metal boundary,'' which is the natural design
variable for a driven, PML-terminated \emph{open radiator}.

We take the differentiable-simulation lens to that underserved corner. The
substrate is a tensor-compiled, differentiable-by-construction 3-D N\'ed\'elec
H(curl) FEM on \emph{unstructured} tetrahedral meshes; the payload is a
moving-boundary shape-adjoint geometry class that treats the positions of the
conductor boundary nodes as the design variables and moves them --- the metal
boundary is represented exactly, at every step, on a body-fitted mesh. The
headline is a design result: end-to-end gradient design of a freeform
curved-metal conformal antenna that drives the entire design band below
$-10$~dB return loss.

We are deliberately narrow about the novelty. Node-motion EM shape adjoints on
unstructured meshes already exist (Section~\ref{sec:related}); this is
explicitly \emph{not} a claim of ``first shape adjoint in EM.'' The contribution
is a specific, previously-unoccupied \emph{combination}: a 3-D frequency-domain
N\'ed\'elec shape adjoint through the \emph{full} open-radiator forward
(box-UPML tensor material $+$ lumped ports $+$ lossy complex-$\varepsilon$ $+$ a
passive $|S_{11}|$ constraint), differentiable-by-construction with
single-factorization reuse, applied to a \emph{freeform, high-DOF} curved-metal
conformal design.

Our contributions are:
\begin{itemize}
  \item A differentiable-by-construction open-radiator FEM \emph{shape} adjoint:
  forward-mode AD Dual-twin element kernels yield exact
  $\partial K/\partial X$, $\partial M/\partial X$, $\partial b/\partial X$
  through the radiating complex-symmetric pencil
  $A(\omega)=K(\nu)-\omega^2 M(\varepsilon)+(j\omega/Z_s)S_p$, so the transpose
  adjoint reuses the \emph{single} forward factorization per band frequency
  ($n_{\text{fact}}=1$), with enforced passivity $|S_{11}|\le 1$ and a
  non-inversion / mesh-distortion guard (Section~\ref{sec:method}).
  \item A high-DOF freeform boundary parametrization (harmonic / P1-Laplace
  mesh-morph, 73 active boundary DOFs) that generalizes single-parameter tuning
  to freeform shape while keeping the volumetric mesh valid, FD-validated to
  relative error $\sim\!10^{-10}$ (Section~\ref{sec:method}).
  \item A headline result: end-to-end gradient design of a curved-metal
  conformal radiator that drives the entire band below $-10$~dB return loss
  (worst-of-band $-5.51 \to -12.06$~dB) on a geometry a Yee-grid density method
  staircases (Section~\ref{sec:results}).
  \item Honest positioning of the narrow white space relative to prior EM
  shape-adjoint and differentiable-FEM work, and an explicit
  \emph{planned}-evaluation protocol for the head-to-head comparisons we have
  not yet run (Sections~\ref{sec:related},~\ref{sec:planned}).
\end{itemize}

\section{Related Work}
\label{sec:related}

\paragraph{Structured-grid density-based differentiable EM (the contrast class).}
The dominant differentiable-EM tools optimize a permittivity/density field on a
Yee/Cartesian grid: Meep's adjoint solver, \textsc{ceviche} and its
challenge suite, SPINS-B~\cite{su_spinsb}, the JAX-based 3-D FDTD framework
FDTDX~\cite{fdtdx}, and the invrs-gym benchmark toolkit~\cite{invrsgym}. These
methods perform topology optimization of a continuous density that is later
binarized. Their canonical successes --- the Piggott et al.\ WDM
demultiplexer~\cite{piggott2015} and the Chung--Miller achromatic
metalens~\cite{chung2020} --- are guided or periodic photonics, not open
radiators terminated by PML. Two references we would like to cite directly ---
the \textsc{ceviche} forward-mode-differentiation-of-Maxwell paper (Hughes et
al., ACS Photonics 2019) and the Meep adjoint-solver documentation --- are not
yet in our verified bibliography and are flagged as citation gaps
(Section~\ref{sec:planned}); we name them here for the reader but do not cite an
unresolved key.

\paragraph{AD shape optimization still on a rectilinear grid.}
Hooten et al.~\cite{hooten2025} accelerate shape optimization with automatic
differentiation, but map shape parameters to a permittivity distribution
\emph{painted on a fixed rectilinear grid}. This reinforces, rather than closes,
the white space we target: the design variable is boundary shape, but the
discretization remains structured and the metal boundary is still resolved by a
density field rather than a body-fitted mesh.

\paragraph{True EM shape (node-motion) adjoints on unstructured meshes exist.}
We are explicit that node-motion EM shape adjoints are not new. Wang and
Anderson~\cite{wang2011} present a discrete Maxwell shape adjoint on
unstructured meshes, but in 2-D, time-domain, lossless, PEC-walled, with
\emph{no PML, no ports, and no S-parameters}. Ghassemi et
al.~\cite{ghassemi2013} evolve a microstrip \emph{antenna} geometry via an FEM
adjoint sensitivity, but with a handful of \emph{low-DOF control vertices} and a
\emph{hand-derived} adjoint. Ham, Mitusch, and Schmidt~\cite{ham2020} provide
\emph{automatic} FEM shape derivatives (Firedrake pyadjoint) for transient PDEs
in general, but not the assembled open-radiator EM forward.

\paragraph{Defensible white space, stated precisely.}
No found work occupies this exact combination: a 3-D frequency-domain N\'ed\'elec
H(curl) FEM shape adjoint through the \emph{full} open-radiator forward
(box-UPML tensor material $+$ lumped ports $+$ lossy complex-$\varepsilon$ $+$
passive $|S_{11}|$), differentiable-by-construction with single-factorization
reuse, applied to a \emph{freeform, high-DOF curved-metal conformal} design. The
three stacked differentiators versus the prior art are: (1) freeform / high-DOF
versus a few control points; (2) full open-radiator physics versus lossless
PEC; and (3) differentiable-by-construction and composable versus a hand-derived
adjoint.

\section{Method}
\label{sec:method}

\subsection{The open-radiator forward pencil}
\label{sec:pencil}

The forward problem is the driven, frequency-domain Maxwell curl-curl equation
discretized with first-order N\'ed\'elec (edge) elements on an unstructured
tetrahedral mesh. At natural frequency $\omega$ the assembled linear system is
the complex-symmetric pencil
\begin{equation}
  A(\omega)\,\mathbf{x} \;=\; \Big[\,K(\nu)\;-\;\omega^2 M(\varepsilon)\;+\;\tfrac{j\omega}{Z_s}\,S_p\,\Big]\,\mathbf{x}
  \;=\; \mathbf{b},
  \label{eq:pencil}
\end{equation}
where $K(\nu)$ is the curl-curl stiffness with reluctivity $\nu=\Lambda^{-1}$,
$M(\varepsilon)$ is the mass matrix with permittivity $\varepsilon=\varepsilon_r\Lambda$,
and $S_p$ is the lumped-port surface term with feed reference impedance
$Z_s$. The radiating boundary is a \emph{matched box-UPML}: an anisotropic
tensor-material absorber whose per-tet real weights $\Lambda$ enter both $K$ and
$M$. Losses in the FR-4 substrate enter as a complex $\varepsilon_r$, making the
pencil genuinely complex-symmetric rather than merely real-shifted. All
frequencies in this work are dimensionless \emph{natural units} as recorded by
the solver; we do not assign physical GHz values.

The scattering parameter is extracted from the port voltage and current,
\begin{equation}
  S_{11} \;=\; \frac{Z-Z_0}{Z+Z_0}, \qquad Z=\frac{V}{I}, \qquad
  I=\frac{2V_{\text{inc}}-V}{R}, \qquad Z_0=R,
  \label{eq:s11}
\end{equation}
with port resistance $R=50$ (dimensionless). Passivity requires
$|S_{11}|\le 1$; we assert this at every forward, gradient, and line-search
evaluation as a physics tripwire.

\subsection{Dual-twin element kernels and the single-factorization shape adjoint}
\label{sec:adjoint}

The design variables are the physical coordinates $X$ of the conductor boundary
nodes. The shape derivative of the objective therefore requires exact
$\partial K/\partial X$, $\partial M/\partial X$, and $\partial b/\partial X$
--- the element-kernel Jacobians with respect to node motion. We obtain these by
\emph{forward-mode automatic differentiation of the closed-form element
kernels}: each entry of the N\'ed\'elec local curl-curl, mass, and current-RHS
kernels is re-evaluated in first-order dual arithmetic $r+d\,\epsilon$
($\epsilon^2=0$). Seeding one node coordinate's tangent with $d=1$ (all others
$0$) returns, in the dual parts, the exact partials of every local matrix entry
with respect to that coordinate. These ``Dual-twin'' kernels mirror the real
$f_{64}$ element kernels entry-for-entry (so their real parts reproduce the
forward assembly), and are shared with the eigenvalue-sensitivity path that
reuses the same element JVP. The tensor-material (box-UPML) twin carries the
fixed real $3\times 3$ per-tet weights $\Lambda$ as lifted constants (zero
tangent), so the tangent flows only through geometry --- the ``fixed $\Lambda$''
convention of the box-UPML shape adjoint.

Because the pencil~\eqref{eq:pencil} is complex-symmetric, the transpose adjoint
system uses the \emph{same} operator as the forward solve. We therefore factor
$A(\omega)$ \emph{once} per band frequency (a single complex LU) and reuse that
factorization for both the forward solve and the adjoint solve:
$n_{\text{fact}}=1$ per frequency. The nodal shape gradient at each band
frequency is $\lambda^{\!\top}(\partial A/\partial X)\,\mathbf{x}$ assembled from
the Dual-twin element Jacobians, where $\lambda$ is the adjoint solution.

\subsection{High-DOF freeform boundary parametrization}
\label{sec:morph}

To move from single-parameter tuning to freeform shape, we parametrize the
boundary with a harmonic mesh-morph. Each free patch-conductor node is an
independent radial-normal design DOF; the interior is extended by a P1-Laplace
mesh-morph regularizer so the volumetric tetrahedra deform gracefully and stay
non-inverted under large boundary deformation. The PML shell, the pinned feed
port, the ground plane, and the outer boundary are held fixed
(\texttt{fixed\_zero}). The multi-DOF design gradient is the per-column
contraction $\partial g/\partial X_p = \langle \text{grad}_{\text{node}}, D_p\rangle$
of the nodal adjoint gradient against the morph columns $D_p$. On the reported
fixture this yields $73$ active freeform boundary DOFs.

The band objective is the frequency-weighted sum of squared reflection
magnitudes,
\begin{equation}
  G(X) \;=\; \sum_{f} w_f\,\big|S_{11}(f;X)\big|^2 ,
  \label{eq:objective}
\end{equation}
with uniform band weights over $\omega\in\{0.30,0.35,0.40\}$ (the v1 spec:
impedance match $+$ bandwidth only). The band nodal gradients are
frequency-weighted, summed, and contracted through the morph columns in one
shot.

\subsection{Guards and validation}
\label{sec:guards}

Two guards run at every evaluation. \emph{Passivity}: $|S_{11}|\le 1$ across the
whole band. \emph{Non-inversion}: a per-tet signed-volume ratio
$|\det(X)/\det(0)|\ge \texttt{min\_vol\_ratio}$ (budget $0.25$) rejects
near-inverted tetrahedra produced by the morph. The full multi-DOF,
multi-frequency band gradient is FD-validated before descent begins: the
analytic directional derivative $\nabla G\cdot\hat d$ is compared against a
central finite difference of the \emph{entire} port $+$ UPML $+$ lossy pipeline
through the public \texttt{driven\_solve\_with\_ports(MatchedUpml)} path, on the
curved mesh, along the steepest-descent direction at $X=0$.

Figure~\ref{fig:setup} sketches the curved conformal radiator inside the
box-UPML with the pinned feed port.

\begin{figure}[t]
  \centering
  \includegraphics[width=0.72\linewidth]{figures/setup_schematic}
  \caption{Schematic of the freeform curved-metal conformal radiator: an FR-4 +
  PEC patch/ground slab wrapped around a cylinder of bend radius $40$ (natural
  units, as recorded) inside a matched box-UPML shell, driven by a pinned-feed
  lumped port on the $x=0$ plane. The free patch-conductor nodes (73 active
  DOFs) move along their radial normals; the PML shell, port plane, ground
  plane, and outer boundary are held fixed. (Placeholder --- rendered by
  \texttt{paper-figures} from \texttt{figures/src/}.)}
  \label{fig:setup}
\end{figure}

\section{Results}
\label{sec:results}

\paragraph{Fixture.}
The reported design is the \texttt{bent\_conformal} fixture: a flat FR-4 + PEC
patch/ground slab wrapped around a cylinder of bend radius $40$ (natural units)
about the $y$-axis to produce genuinely curved conformal metal, with an $8$
(natural units) box-UPML shell and an $x=0$ port plane. The discretization has
$6173$ edges, $997$ nodes, and $4875$ tetrahedra; the design vector has $73$
active freeform boundary DOFs and a $50$ (dimensionless) port resistance.

\paragraph{Gradient correctness.}
The full band gradient matches a central finite difference of the public driven
pipeline on the curved mesh to relative error $1.17\times10^{-10}$ (analytic and
central directional derivatives both $1.751288204\times10^{-1}$, FD step
$10^{-6}$, tolerance $5\times10^{-3}$), with one factorization per band
frequency.

\paragraph{Headline design result.}
Bounded backtracking steepest descent (Armijo line search) on the 73-DOF design
vector reaches a genuine \texttt{target\_reached} terminal in $6$ accepted
steps (of a $600$-step cap --- a true target hit, not a step cap), with $0$
total backtracks. The band objective~\eqref{eq:objective} falls from
$1.897943102\times10^{-1}$ to $3.411688781\times10^{-2}$, a reduction factor of
$5.56$. The worst-of-band return loss improves from $-5.51$~dB to $-12.06$~dB,
driving the \emph{entire} design band below the $-10$~dB spec. The final
per-frequency reflection is
\begin{center}
\begin{tabular}{lccc}
\toprule
$\omega$ & $0.30$ & $0.35$ & $0.40$ \\
\midrule
$|S_{11}|$ (dB) & $-12.06$ & $-23.92$ & $-14.42$ \\
$|S_{11}|$ (mag) & $0.2494$ & $0.0637$ & $0.1900$ \\
\bottomrule
\end{tabular}
\end{center}
Figure~\ref{fig:s11band} shows the convergence trajectory and the final
per-frequency return loss against the $-10$~dB target.

\begin{figure}[t]
  \centering
  \includegraphics[width=0.9\linewidth]{figures/s11_band}
  \caption{Left: worst-of-band return loss versus optimization step
  ($-5.51 \to -12.06$~dB over $6$ accepted steps). Right: final per-frequency
  $|S_{11}|$ (dB) at $\omega\in\{0.30,0.35,0.40\}$ against the $-10$~dB spec
  (dashed). All values from the committed artifact
  \texttt{conformal\_results.toml}. (Placeholder --- rendered by
  \texttt{paper-figures} from \texttt{figures/src/plot\_s11\_band.py}.)}
  \label{fig:s11band}
\end{figure}

\paragraph{Guards and determinism.}
The non-inversion guard never bound: the worst per-tet signed-volume ratio was
$0.572$ against the $0.25$ budget, at a maximum nodal displacement of $0.560$
(natural units). Passivity $|S_{11}|\le 1$ held at every evaluation. The
gradient norm fell from $1.751\times10^{-1}$ to $8.987\times10^{-2}$. An
independent fresh forward band solve at the final design (public
\texttt{driven\_solve\_with\_ports}, no adjoint) reproduces the optimizer's
final band objective to relative $0.0$, and the run is bit-identical on re-run.

\paragraph{Methodological foundation.}
The composed driven SHAPE-adjoint stack --- lossy complex-$\varepsilon$ adjoint,
pinned/moving-feed lumped-port adjoint, box-UPML tensor-material adjoint --- was
validated by a prior single-DOF capstone that retuned a flat patch, at
FD relative error $\sim\!10^{-9}$. The conformal result extends that from $1$
DOF to $73$ freeform DOFs on curved geometry.

\section{Discussion}
\label{sec:discussion}

The result substantiates a narrow but clean claim: a differentiable-by-construction
open-radiator FEM can \emph{reach} a freeform curved-conformal $-10$~dB design
by moving conductor-boundary nodes on a body-fitted mesh, with an exact,
single-factorization shape adjoint and enforced physics guards. The design space
is genuinely high-dimensional (73 independent DOFs; not one parametric knob),
and the geometry is genuinely curved --- the two properties that motivate a
body-fitted node-motion adjoint over a structured density field.

\paragraph{Limitations and threats to validity.}
The scope of v1 is impedance match $+$ bandwidth only; radiation pattern and
gain (a near-to-far-field adjoint) are not addressed here
(Section~\ref{sec:planned}). The reported design is a single fixture at one bend
radius and one band; we have not swept mesh resolution, bend radius, or band
placement, so we make no claim about generality across radiator families. The
comparative superiority over structured-grid density methods is
\emph{not} yet quantified --- we claim reachability, not a measured win
(Section~\ref{sec:planned}). We also note a reproducibility detail in the spirit
of full disclosure: an initial optimization run mis-recorded a premature
(step-cap) stop as a physics limit; a corrected line search reached the true
\texttt{target\_reached} terminal, which is the run reported here.

\section{Planned Evaluation (Future Work)}
\label{sec:planned}

The following evaluations are \emph{planned and not yet run}; we report
\emph{no} baseline numbers for them and make no comparative claim beyond
reachability.

\paragraph{Three-way head-to-head.}
The intended comparison pits the present body-fitted node-motion adjoint against
(i) an FDTD/FDFD-density optimizer (e.g.\ Meep-adjoint or \textsc{ceviche})
forced to staircase the curved conformal geometry, and (ii) a low-DOF
parametric antenna sweep. The metric is match quality at fixed compute /
resolution budget. This is the announced next step; the honest claim today is
that GEODE reaches a freeform curved-conformal $-10$~dB design, while the
quantified superiority over FDTD-density remains to be measured.

\paragraph{Radiation pattern / gain.}
A near-to-far-field (NTFF) shape adjoint extending the objective from
$|S_{11}|$ to pattern/gain is future work and is not claimed here.

\paragraph{Citation gaps.}
Three references are named in the text but are \emph{not} yet in our
resolver-verified bibliography, and are surfaced here as gaps for a dedicated
literature-search pass rather than cited from an unresolved key: (1) the
\textsc{ceviche} forward-mode-differentiation-of-Maxwell paper (Hughes et al.,
ACS Photonics 2019) --- needs a resolvable DOI/arXiv id; (2) the Meep
adjoint-solver reference / documentation; and (3) a canonical antenna
shape/topology-optimization reference beyond Ghassemi et al.\ that would
strengthen the parametric-baseline framing.

\section{Conclusion}
\label{sec:conclusion}

We presented a tensor-compiled, differentiable-by-construction 3-D N\'ed\'elec
H(curl) FEM that exposes an exact, single-factorization discrete shape adjoint
through the full open-radiator forward --- matched box-UPML tensor material,
lumped ports, lossy complex-$\varepsilon$, and a passive $|S_{11}|$ --- on
unstructured tetrahedral meshes. Driving a 73-DOF harmonic freeform boundary
morph, we gradient-designed a curved-metal conformal radiator that drives the
entire design band below $-10$~dB return loss (worst-of-band $-5.51 \to
-12.06$~dB), with the band gradient FD-validated to relative error
$1.17\times10^{-10}$, enforced passivity, and a non-inversion guard. The
contribution is a narrow, previously-unoccupied combination --- freeform /
high-DOF, full open-radiator physics, differentiable-by-construction with
single-factorization reuse --- explicitly not a first EM shape adjoint. The
quantified head-to-head against structured-grid density methods, and a
pattern/gain NTFF adjoint, are the announced next steps.

\bibliographystyle{plainnat}
\bibliography{refs}

\end{document}
